%% -----------------------------------------------------------------------------------------------
\section{Physical-record measurements}\label{sec:s-physics}

Three data sources are used and kept apart: (A) the original development record (C1, driver 550.120); (B) a static two-box experiment; (C) a sub-step re-run of 11 used development scene-runs on the calibration host (driver 595.99.02). Only (A) is the record analysed in the main text; (B) and (C) are post hoc and explain part of a failure found in (A).

\paragraph{Definitions (source A)} The checks are implemented in \path{src/hci/spatial_links.py}. Logged poses are taken at the end of each control step. For every listed PhysX contact pair, each contact point carries a signed separation and an impulse; a pair with any non-zero impulse is an \emph{observed contact}, a listed pair with zero impulse is \emph{proximity}. The reconstructed signed distance of two shapes is computed from the logged poses and collision shapes. \emph{S4}: during the move (control steps after the trigger), the minimum reconstructed distance between any robot shape and the target cube is positive. \emph{S5}: every observed-contact shape pair in any phase of the branch is reconstructed at most 1\,mm apart; overlaps are reported (maximum and median penetration) without a threshold. \emph{S6}: every listed shape pair whose PhysX separation is positive, in any phase, has $|\text{PhysX separation}-\text{reconstructed distance}|\le1$\,mm. A row of S6 is one such pair at one control step.

\paragraph{S4 and S5 (source A)} S4 failed only in the oracle branch of scene 51152 (seed 2). In that scene the trigger fired while the fingers held the target cube B: from the steps before the trigger onwards the reconstructed finger--target distance was $-2.16$\,mm with about 27.6\,N of contact force on each finger, cube A was not grasped, and the oracle's 12-step move pushed cube B by 23.6\,mm and enlarged the offset from 52 to 76\,mm (C0 failed as well, so this was not a harm). All vision arms abstained in this scene for lack of a held-cube estimate. S5 failed in 11 of 3,874 branch records; these share 2 distinct events with fast-moving cubes (112--189\,mm/s), in scene 51143 (seed 2) and scene 51270 (seed 1), in cube--cube and finger--cube pairs.

\paragraph{S6 (source A)} 54,267 of 580,539 rows (9.3\%) exceed 1\,mm; per pair kind see Fig.~7a of the main text. A second, independent reconstruction on a subset (8 to about 1,000 rows per pair kind) agreed with the first to within $10^{-12}$\,mm for pairs of box-shaped geometry (cubes, fingers, camera, table), within $4\times10^{-6}$\,mm for cube B--table, and within 0.22--0.24\,mm for hand--cube pairs, where the hand is a triangle mesh; the S6 disagreement is therefore between PhysX and the logged poses, not between the two reconstructions.

\paragraph{Static known geometry (source B)} 59 static cases (aligned, rotated, edge and vertex configurations of two boxes; gaps from $-3$ to $40$\,mm) were simulated for one control step. In the 56 cases with listed contact points, the PhysX separation equals the analytic gap \emph{before} the last sub-step within 0.05\,mm.

\paragraph{Sub-step re-run (source C)} Logging every one of the five physics sub-steps per control step on 11 used scene-runs gave 56,558 separated listed pairs at the last sub-step. Compared with the pose after the step, 5,181 exceed 1\,mm (median absolute difference 0.044\,mm); compared with the pose before the last sub-step, 1,874 (median $2\times10^{-5}$\,mm). By relative speed: below 1\,mm/s 137 $\to$ 137; 1--10\,mm/s 690 $\to$ 681; 10--100\,mm/s 918 $\to$ 910; above 100\,mm/s 3,436 $\to$ 146. The residual 1,874 rows are zero-impulse proximity pairs of a finger box and a cube in which PhysX reports a larger separation than the reconstruction. Shape-index mismatch (17 of 45,645) and static edge geometry (20 dedicated cases, zero difference to a GJK reference) were excluded. The traces of this re-run differ from the original development traces in all 11 scene-runs (with different outcomes in 6), so source C describes the simulator on the calibration host, not the original record.

\paragraph{Hull sensitivity (source A, post hoc)} Replacing each logged link mesh by the convex hull of its vertices (volume excess at most 0.29\%) and recomputing S4 and S5 on all 3,874 branch records changed no status. Hand--target distances changed by at most 0.24\,mm and arm-link distances by at most 0.03\,mm; finger distances, which determine the S4 minimum in these branches, did not change.
